要掌握主存储器与CPU的连接之前，首先需要了解一个考研必考的知识点，即\textbf{{存储器容量的扩充。}}{一般存储器容量的扩充通常有3类：}\textbf{位扩充、字扩充和字位扩充}{。}

\textbf{1.
位扩充（增加{a}\includegraphics[width=0.09375in,height=0.08333in]{texmath/ff0c81times}{b}后面的b）}

位扩充指增加存储字长，例如，现要将1K\textbf{\includegraphics[width=0.09375in,height=0.08333in]{texmath/ff0c81times}}4位的芯片组成1K\textbf{\includegraphics[width=0.09375in,height=0.08333in]{texmath/ff0c81times}}8位的存储器。

\textbf{2.
字扩充（增加{a}\includegraphics[width=0.09375in,height=0.08333in]{texmath/ff0c81times}{b}前面的a）}

字扩充是增加存储单元的个数，例如，现要将1K\textbf{\includegraphics[width=0.09375in,height=0.08333in]{texmath/ff0c81times}}8位的芯片组成2K\textbf{\includegraphics[width=0.09375in,height=0.08333in]{texmath/ff0c81times}}8位的存储器。

\textbf{3.
字位扩充（增加{a}\includegraphics[width=0.09375in,height=0.08333in]{texmath/ff0c81times}{b}中的a和b）}

字位扩充既增加了存储单元的个数，又增加了存储字长，例如，要用1K\textbf{\includegraphics[width=0.09375in,height=0.08333in]{texmath/ff0c81times}}4位的芯片组成4K\textbf{\includegraphics[width=0.09375in,height=0.08333in]{texmath/ff0c81times}}8位的存储器。

{\textbf{注意：在进行字位扩充时，一定是先进行位扩充，再进行字扩充。}}

{ \textbf{} }

\textbf{{\textbf{另外，记住下面这个芯片数量计算公式：}}}

\textbf{}

{如果要求将容量为a}\includegraphics[width=0.09375in,height=0.08333in]{texmath/ff0c81times}{b的芯片组成容量为c}\includegraphics[width=0.09375in,height=0.08333in]{texmath/ff0c81times}{d的芯片，假设需要芯片的数量为n，则n=(c*d)/(a*b)（该公式就是整个存储器的容量除以单个芯片的容量）}。

以上掌握之后，只需要根据需要组成的芯片计算地址线和数据线的数量，再使用译码器（根据芯片的数量来选择，一般有2-4译码器和3-8译码器）连接起来即可。
